% Generated draft: requires mathematical and source review.
\documentclass[11pt]{article}
\usepackage[margin=25mm]{geometry}
\usepackage{fontspec}
\setmainfont{Latin Modern Roman}
\usepackage{amsmath,amssymb,mathtools,unicode-math}
\setmathfont{Latin Modern Math}
\usepackage{enumitem,needspace,xurl,hyperref}
\setlength{\parindent}{0pt}
\setlength{\parskip}{5pt}
\setlength{\emergencystretch}{3em}
\newcommand{\N}{\mathbb{N}}
\newcommand{\Z}{\mathbb{Z}}
\newcommand{\Q}{\mathbb{Q}}
\newcommand{\R}{\mathbb{R}}
\newcommand{\C}{\mathbb{C}}
\newcommand{\F}{\mathbb{F}}
\newcommand{\A}{\mathbb{A}}
\newcommand{\PP}{\mathbb{P}}
\newcommand{\E}{\mathbb{E}}
\newcommand{\eps}{\varepsilon}
\newcommand{\dd}{\,\mathrm{d}}
\begin{document}
\section*{UnsolvedMath ID: 30001319}
\textbf{Source problem code:} catalogue entry
\subsection{Definitions and mathematical statement}
The supplied question concerns arm events in two-dimensional critical percolation. It does not specify the lattice, site or bond percolation, boundary conditions, or the precise arm-event convention. The following gives a conditional annulus formulation; these choices must be fixed before it specifies a unique problem.

Fix a planar percolation model at its critical parameter, with probability law \(\mathbb P_c\). For each admissible positive integer \(j\), choose a fixed inner radius \(r_j>0\) for which the relevant \(j\)-arm events are possible. For \(n>r_j\), let \(A_j^{\mathrm{mono}}(r_j,n)\) denote the event that the annular region between radii \(r_j\) and \(n\) contains \(j\) mutually disjoint arms connecting its inner and outer boundaries, all of one specified colour. The interpretation of colour, the underlying paths, and disjointness must follow the chosen model; in particular, closed arms may require a dual or matching-lattice convention. The source does not specify which monochromatic colour is intended.

For \(j\geq 2\), let \(\sigma=(\sigma_1,\ldots,\sigma_j)\in\{\mathrm{open},\mathrm{closed}\}^j\) contain both colours, and let \(A_{j,\sigma}^{\mathrm{poly}}(r_j,n)\) denote the corresponding event with arm colours prescribed by \(\sigma\) in cyclic order. When it exists, define its polychromatic exponent by
\[
\alpha_{j,\sigma}
=-\lim_{n\to\infty}
\frac{\log\mathbb P_c\!\left(A_{j,\sigma}^{\mathrm{poly}}(r_j,n)\right)}{\log n}.
\]
No independence from the colour sequence, the inner radius, or the model is assumed here.

The source asks whether, for each intended admissible \(j\), there is a finite exponent \(\alpha'_j\geq 0\) such that
\[
\mathbb P_c\!\left(A_j^{\mathrm{mono}}(r_j,n)\right)
=n^{-\alpha'_j+o(1)}
\qquad(n\to\infty),
\]
where \(o(1)\) tends to zero with the model, \(j\), colour, and inner radius fixed. Equivalently, it asks whether the finite limit
\[
\alpha'_j
=-\lim_{n\to\infty}
\frac{\log\mathbb P_c\!\left(A_j^{\mathrm{mono}}(r_j,n)\right)}{\log n}
\]
exists. It further asks for the relation between these monochromatic exponents and the corresponding polychromatic arm exponents. The supplied record gives neither a proposed identity nor a specific comparison to be established, and it does not state the intended range of \(j\) or the scope of models covered.

The supplied dated curation labels the problem partially solved: it reports existence and strict bounds for monochromatic arm exponents, a later exact evaluation of the monochromatic two-arm exponent, and remaining questions about exact values for the full higher-arm family and the general relation. This is a supplied literature assessment, not an independently verified theorem statement; its model-specific hypotheses and exact conclusions are not included in the record.
\subsection{Short English statement}
For a specified two-dimensional critical percolation model, ask whether the probability of \(j\) disjoint arms of one colour reaching distance \(n\) has a well-defined power-law exponent, and how that exponent compares with exponents for arms containing both colours. The model, arm conventions, and intended comparison need clarification. Supplied dated curation describes partial progress; this statement-only draft requires human mathematical review.
\subsection{Sources}
\begin{itemize}
\item Supplied local UnsolvedMath registry record, permanent ID \(30001319\), source problem code \texttt{catalogue entry}, title \emph{Monochromatic Arm Exponents in Critical Percolation}, in the Oberwolfach Reports Open Problems collection. No individual proposer is supplied. The registry records proposed year \(2009\).
\item Supplied report citation: \emph{Scaling Limits in Models of Statistical Mechanics} (2010), Oberwolfach Reports, pp.\ 2115--2178, DOI \texttt{10.4171/owr/2009/38}. The literature-source label is ``original Oberwolfach report 4081'': \url{https://doi.org/10.4171/owr/2009/38}. The record describes Oberwolfach Reports as documenting mathematical workshops and their open-problem discussions. No precise page or passage for this question is supplied.
\item Supplied literature reference labelled ``Beffara-Nolin theorem'': \url{https://arxiv.org/abs/0906.3570}. The record associates this reference with existence and strict bounds for monochromatic arm exponents, without supplying the precise theorem or hypotheses.
\item Supplied literature reference labelled ``Nolin-Qian-Sun-Zhuang backbone exponent'': \url{https://arxiv.org/abs/2309.05050}. The record associates this reference with an exact evaluation of the monochromatic two-arm exponent, without supplying the precise theorem or hypotheses.
\item Literature assessment dated exactly \texttt{2026-08-21}: the supplied curation labels the problem partially solved, reports existence and strict bounds for monochromatic arm exponents and an exact monochromatic two-arm exponent, and describes exact values for the full higher-arm family and the general relation as remaining open. This is supplied dated curation derived from the OWR record, not independently verified. The source explicitly requests independent verification of time-sensitive status and cited claims.
\item The registry labels the statement \texttt{corrected\_verified}, with the supplied explanation that the arm-exponent questions are complete and that indices, exponent notation, and asymptotic notation were restored. This records the registry's verification claim only; no source was independently visited for this draft.
\end{itemize}
\end{document}
